
\section{Introduction}

\label{sec:intro}


Open-domain question answering (QA)~\cite{voorhees1999trec} is a task that answers factoid questions using a large collection of documents.
While early QA systems are often complicated and consist of multiple components~(\citet{ferrucci2012introduction,moldovan2003performance}, \textit{inter alia}),
the advances of reading comprehension models suggest a much simplified two-stage framework: (1) a context \emph{retriever} first selects a small subset of passages where some of them contain the answer to the question, and then (2) a machine \emph{reader} can thoroughly examine the retrieved contexts and identify the correct answer~\cite{chen2017reading}.
Although reducing open-domain QA to machine reading is a very reasonable strategy, a huge performance degradation is often observed in practice\footnote{For instance,
the exact match score on SQuAD v1.1 drops from above 80\% to less than 40\%~\cite{yang2019end}.}, indicating the needs of improving retrieval.


Retrieval in open-domain QA %
is usually implemented using TF-IDF or BM25~\cite{robertson2009probabilistic}, which matches keywords efficiently with an inverted index and can be seen
as representing the question and context in high-dimensional, sparse vectors (with weighting).
Conversely, the \emph{dense}, latent semantic encoding is \emph{complementary} to sparse representations by design.
For example, synonyms or paraphrases
that consist of completely different tokens may still be mapped to vectors close to each other.
Consider the question \ti{``Who is the bad guy in lord of the rings?''}, which can be answered from the context
\ti{``Sala Baker is best known for portraying the villain Sauron in the Lord of the Rings trilogy.''}
A term-based system would have difficulty retrieving such a context, while a dense retrieval system would be able to better match \ti{``bad guy''} with \ti{``villain''}
and fetch the correct context.
Dense encodings are also \emph{learnable} by adjusting the embedding functions, which provides additional flexibility to have a task-specific representation. With special in-memory data structures and indexing schemes, retrieval can be done efficiently using maximum inner product search (MIPS) algorithms~(e.g.,~\citet{NIPS2014_5329, guo2016quantization}).




However, it is generally believed that learning a good dense vector representation needs a large number of labeled pairs of question and contexts.
Dense retrieval methods have thus never be shown to outperform TF-IDF/BM25 for open-domain QA before ORQA~\cite{lee2019latent}, which proposes a sophisticated inverse cloze task (ICT) objective, predicting the blocks that contain the masked sentence, for additional pretraining.
The question encoder and the reader model are then fine-tuned using pairs of questions and answers jointly. %
Although ORQA successfully demonstrates that dense retrieval can outperform BM25, setting new state-of-the-art results on multiple open-domain QA datasets,
it also suffers from two weaknesses.
First, ICT pretraining is computationally intensive and it is not completely clear that regular sentences are good surrogates of questions in the objective function.
Second, because the context encoder is not fine-tuned using pairs of questions and answers, the corresponding representations could be suboptimal.


In this paper, we address the question: can we train a better dense embedding model using only pairs of questions and passages (or answers), \emph{without} additional pretraining?
By leveraging the now standard BERT pretrained model~\cite{devlin2018bert} and a dual-encoder architecture~\cite{bromley1994signature}, we focus on developing the right training scheme using a relatively small number of question and passage pairs. %
Through a series of careful ablation studies, our final solution is surprisingly simple: the embedding is optimized for maximizing inner products of the question and relevant passage vectors, with an objective comparing all pairs of questions and passages in a batch.
Our \emph{Dense Passage Retriever} (DPR) is exceptionally strong. It not only outperforms BM25 by a large margin (65.2\% vs. 42.9\% in Top-5 accuracy), but also results in a substantial improvement on the end-to-end QA accuracy compared to ORQA (41.5\% vs. 33.3\%) in the open Natural Questions setting~\cite{lee2019latent,kwiatkowski2019natural}.


Our contributions are twofold.
First, we demonstrate that with the proper training setup, simply fine-tuning the question and passage encoders on existing question-passage pairs
is sufficient to greatly outperform BM25.
Our empirical results also suggest that additional pretraining may not be needed.
Second, we verify that, in the context of open-domain question answering, a higher retrieval precision indeed translates to a higher end-to-end QA accuracy.
By applying a modern reader model to the top retrieved passages, we achieve comparable or better results on multiple QA datasets in the open-retrieval setting,
compared to several, much complicated systems.


